% 预算 B=7 下两种逐轮规则（传感器覆盖实例，成本 a4 b3 c4 d1）。
% 总增益规则：第1轮选a（+10，剩3），第2轮可负担 b(+2)、d(+1)，选b，剩0，覆盖12。
% 单位成本规则：第1轮比值 a2.5 b2.67 c2 d1，选b（+8，剩4）；第2轮 a 1、c 2、d 0，选c（+8），剩0，覆盖16（枚举最优）。
\begin{tikzpicture}[x=8mm,y=1mm,line width=.6pt,font=\figurefont\fontsize{8}{10}\selectfont]
  \foreach \yy/\ttl in {11/总增益最大,0/单位成本增益最大}{
    \node[FigureInk,anchor=east] at (-0.2,\yy+3) {\ttl};
    \draw[FigureGrid] (0,\yy) rectangle (7,\yy+6);
    \foreach \t in {1,...,6}{\draw[FigureGrid,line width=.4pt] (\t,\yy)--(\t,\yy+6);}
  }
  % 总增益规则
  \filldraw[fill=FigureBlueFill,draw=FigureBlue,rounded corners=1pt] (0.05,11.5) rectangle (3.95,16.5);
  \node[FigureInk] at (2,14) {$a$：成本4，$+10$};
  \filldraw[fill=FigureBlueFill,draw=FigureBlue,rounded corners=1pt] (4.05,11.5) rectangle (6.95,16.5);
  \node[FigureInk] at (5.5,14) {$b$：成本3，$+2$};
  \node[FigureInk,anchor=west] at (7.2,14) {覆盖$12$};
  % 单位成本规则
  \filldraw[fill=FigureGreenFill,draw=FigureGreen,rounded corners=1pt] (0.05,0.5) rectangle (2.95,5.5);
  \node[FigureInk] at (1.5,3) {$b$：成本3，$+8$};
  \filldraw[fill=FigureGreenFill,draw=FigureGreen,rounded corners=1pt] (3.05,0.5) rectangle (6.95,5.5);
  \node[FigureInk] at (5,3) {$c$：成本4，$+8$};
  \node[FigureInk,anchor=west] at (7.2,3) {覆盖$16$};
  % 预算刻度
  \foreach \t in {0,...,7}{\node[FigureMuted,font=\FigureNoteFont\fontsize{7.5}{9}\selectfont] at (\t,-3) {\t};}
  \node[FigureMuted,anchor=west,font=\FigureNoteFont\fontsize{7.5}{9}\selectfont] at (7.2,-3) {已用预算};
\end{tikzpicture}
